\chapter{Аналитический раздел}

В данном разделе представляются анализ предметной области, понятия сетевого протокола, модели ввода-вывода, виды запросов к серверу, описания механизмов prefork, select, мультиплексирования и требований к безопасности сервера.

\section{Анализ предметной области}

Статический сервер -- это программное обеспечение, предназначенное для предоставления пользователям через сеть статического содержимого по их запросу. Статическое содержимое представляет из себя неизменяемые данные, такие как HTML-страницы, изображения, файлы.

\section{Сетевой протокол}

Для обеспечения передачи данных по организованным каналам связи в инфокоммуникационных системах применяются протоколы. Протокол~--~набор соглашений интерфейса логического уровня, которые определяют обмен данными между различными программами. По сути задача протокола~--~предоставить передаваемые пользователем данные в приемлемой форме и сопроводить служебной информацией, которая позволит корректно передать по каналам связи и интерпретировать данные на принимаемой стороне~\cite{tihomirova}. 

Сетевой протокол~--~набор правил и действий (очерёдности действий), позволяющий осуществлять соединение и обмен данными между двумя и более включёнными в сеть устройствами. 

\subsection{Протокол HTTP}

\section{Модели ввода-вывода}

\subsection{Синхронная модель ввода-вывода}

\section{Запросы к серверу}

\section{Механизм prefork}

\section{Механизм select}

\section{Мультиплексирование}

\section{Безопасность сервера}

\section*{Вывод}


